\documentclass[compress, 8pt, aspectratio=169]{beamer}

\usepackage{presentationtemplate}
\usepackage[fontsize=\small,linenosfontsize=\tiny,askip=3mm]
  {mylisting}
\usepackage[fontsize=\small,askip=3mm]{terminal}
\usepackage{textcomp}
\usepackage[autostyle=true]{csquotes}

\title{Фазы трансляции}

\makeatletter
\newdimen\tp@FVleftmargin
\let\tp@FVSetupFont\FV@SetupFont
\def\FV@SetupFont{%
  \tp@FVleftmargin=\leftmargin
  \tp@FVSetupFont
  \let\leftmargin\tp@FVleftmargin}
\makeatother

\begin{document}

\frame[plain]{\titlepage}

% ------------------------------------------------------------------
% Слайд 2. Фазы трансляции
% ------------------------------------------------------------------

\begin{frame}

    \frametitle{Фазы трансляции}

    \begin{enumerate}

        \item Mapping source characters

        \item Splicing lines

        \item Lexing

        \item \textbf{Preprocessing (препроцессор)}

        \item Converting character constants and string literals
            to execution character set

        \item Concatenating string literals

        \item \textbf{Compiling (компиляция)}

        \item \textbf{Linking (линковка)}

    \end{enumerate}

    \bigbreak

    Более подробно:
    \href{https://en.cppreference.com/w/c/language/translation_phases}
        {eng},
    \href{https://ru.cppreference.com/w/c/language/translation_phases}
        {ru}.

\end{frame}

% ------------------------------------------------------------------
% Слайд 3. Компиляторы C
% ------------------------------------------------------------------

\begin{frame}

    \frametitle{Компиляторы C}

    \begin{itemize}

        \item \texttt{\textbf{gcc}} от GNU Compiler Collection (GCC)

        \item \texttt{\textbf{clang}} от Low Level Virtual Machine (LLVM)

        \item Microsoft Visual C (MSVC)

        \item и др.

    \end{itemize}

\end{frame}

% ------------------------------------------------------------------
% Слайд 4. Препроцессор: схема
% ------------------------------------------------------------------

\begin{frame}

    \frametitle{Препроцессор}

    \centering

    \begin{tikzpicture}[node distance = 5cm]

        \node (source1) {
            \includegraphics[width=1.7cm]
                {presentations/images/source_code.png}
        };
        \node (source2) [right=of source1] {
            \includegraphics[width=1.7cm]
                {presentations/images/source_code.png}
        };

        \draw[->, ultra thick] (source1.east) -- (source2.west)
            node[midway, above, align=center]
                {Работа\\препроцессора};

    \end{tikzpicture}

\end{frame}

% ------------------------------------------------------------------
% Слайд 5. Препроцессор: директива #define
% ------------------------------------------------------------------

\begin{frame}[fragile]

    \frametitle{Препроцессор: директива \\ \texttt{\#define}}

    \begin{columns}[T]

        \begin{column}{0.5\textwidth}

            \myinputlisting[minted language=c]
                {presentations/05-translation-phases/src/}
                {defines-example.c}

        \end{column}

        \begin{column}{0.5\textwidth}

            \begin{terminalwindow}
!\shellcommand{gcc \myhighlight{yellow}{-E} defines-example.c}!
# 0 "defines-example.c"
# 0 "<built-in>"
# 0 "<command-line>"
# 1 "/nix/store/!\dots!/include/stdc-predef.h" 1 3 4
# 0 "<command-line>" 2
# 1 "defines-example.c"




int main(void) {
    int i = 1 + 32;
    return 0;
}
            \end{terminalwindow}

        \end{column}

    \end{columns}

\end{frame}

% ------------------------------------------------------------------
% Слайд 6. Препроцессор: директива #define (макрос с параметром)
% ------------------------------------------------------------------

\begin{frame}[fragile]

    \frametitle{Препроцессор: директива \\ \texttt{\#define}}

    \begin{columns}[T]

        \begin{column}{0.5\textwidth}

            \myinputlisting[minted language=c]
                {presentations/05-translation-phases/src/}
                {defines-param-example.c}

        \end{column}

        \begin{column}{0.5\textwidth}

            \begin{terminalwindow}
!\shellcommand{gcc defines-param-example.c}!
!\shellcommand{./a.out}!
1, 3.14, text
            \end{terminalwindow}

        \end{column}

    \end{columns}

\end{frame}

% ------------------------------------------------------------------
% Слайд 7. Препроцессор: директива #include
% ------------------------------------------------------------------

\begin{frame}[fragile]

    \frametitle{Препроцессор: директива \\ \texttt{\#include}}

    \begin{columns}[T]

        \begin{column}{0.5\textwidth}

            \myinputlisting[minted language=c]
                {presentations/05-translation-phases/src/}
                {include-header.h}

            \myinputlisting[minted language=c]
                {presentations/05-translation-phases/src/}
                {include-source.c}

        \end{column}

        \begin{column}{0.5\textwidth}

            \begin{terminalwindow}
!\shellcommand{gcc -E include-source.c}!
# 0 "include-source.c"
# 0 "<built-in>"
# 0 "<command-line>"
# 1 "/nix/store/!\dots!/include/stdc-predef.h" 1 3 4
# 0 "<command-line>" 2
# 1 "include-source.c"
# 1 "include-header.h" 1
!\myhighlight{yellow}{int a = 0;}!
# 2 "include-source.c" 2

int main(void) {
    int b = a;
    return 0;
}
            \end{terminalwindow}

        \end{column}

    \end{columns}

\end{frame}

% ------------------------------------------------------------------
% Слайд 8. Препроцессор: директива #pragma
% ------------------------------------------------------------------

\begin{frame}[fragile]

    \frametitle{Препроцессор: директива \\ \texttt{\#pragma}}

    \begin{columns}[T]

        \begin{column}{0.5\textwidth}

            \myinputlisting[minted language=c]
                {presentations/05-translation-phases/src/}
                {pragma-example.h}

            \myinputlisting[minted language=c]
                {presentations/05-translation-phases/src/}
                {pragma-example.c}

        \end{column}

        \begin{column}{0.5\textwidth}

            \begin{terminalwindow}
!\shellcommand{gcc -E pragma-example.c}!
# 0 "pragma-example.c"
# 0 "<built-in>"
# 0 "<command-line>"
# 1 "/nix/store/!\dots!/include/stdc-predef.h" 1 3 4
# 0 "<command-line>" 2
# 1 "pragma-example.c"
# 1 "pragma-example.h" 1
       

const double PI = 3.14;
# 2 "pragma-example.c" 2




int main(void) {
    double pi = PI;
}
            \end{terminalwindow}

        \end{column}

    \end{columns}

\end{frame}

% ------------------------------------------------------------------
% Слайд 9. Компиляция: схема
% ------------------------------------------------------------------

\begin{frame}

    \frametitle{Компиляция}

    \centering

    \begin{tikzpicture}[node distance = 5cm]

        \node (source) {
            \includegraphics[width=1.7cm]
                {presentations/images/source_code.png}
        };
        \node (binary) [right=of source] {
            \includegraphics[width=1.7cm]
                {presentations/images/binary-file.png}
        };
        \node [below=0mm of binary] { \texttt{*.o} };

        \draw[->, ultra thick] (source.east) -- (binary.west)
            node[midway, above] {Компиляция};

    \end{tikzpicture}

\end{frame}

% ------------------------------------------------------------------
% Слайд 10. Компиляция: objdump -f и -tC
% ------------------------------------------------------------------

\begin{frame}[fragile]

    \frametitle{Компиляция}

    \begin{columns}[T]

        \begin{column}{0.25\textwidth}

            \myinputlisting[minted language=c,
                minted options={
                    fontsize=\scriptsize,
                    breaklines,
                    breakanywhere,
                    linenos,
                    escapeinside=||,
                    numbersep=2pt
                }]
                {presentations/05-translation-phases/src/}
                {comp-example.c}

        \end{column}

        \begin{column}{0.75\textwidth}

            \begin{terminalwindow}
!\shellcommand{gcc \myhighlight{yellow}{-c} comp-example.c -o comp-example.o}!
!\shellcommand{objdump -f comp-example.o}!

comp-example.o:     file format !\myhighlight{green}{elf64-x86-64}!
architecture: i386:x86-64, flags 0x00000011:
HAS_RELOC, HAS_SYMS
start address 0x0000000000000000
!\shellcommand{objdump \myhighlight{yellow}{-tC} comp-example.o}!

comp-example.o:     file format elf64-x86-64

SYMBOL TABLE:
0000000000000000 l    df *ABS*  0000000000000000 comp-example.c
0000000000000000 l    d  .text  0000000000000000 .text
0000000000000000 l    d  .rodata        0000000000000000 .rodata
0000000000000000 g     F .text  000000000000000e !\myhighlight{green}{PI}!
000000000000000e g     F .text  000000000000001b  !\myhighlight{green}{main}!
            \end{terminalwindow}

        \end{column}

    \end{columns}

\end{frame}

% ------------------------------------------------------------------
% Слайд 11. Компиляция: objdump -d (целиком)
% ------------------------------------------------------------------

\begin{frame}[fragile]

    \frametitle{Компиляция}

            \begin{terminalwindow}
!\shellcommand{objdump \myhighlight{yellow}{-d} comp-example.o}!

comp-example.o:     file format elf64-x86-64


Disassembly of section .text:

0000000000000000 <PI>:
   0:   55                      push   %rbp
   1:   48 89 e5                mov    %rsp,%rbp
   4:   f3 0f 10 05 00 00 00    movss  0x0(%rip),%xmm0        # c <PI+0xc>
   b:   00 
   c:   5d                      pop    %rbp
   d:   c3                      ret

000000000000000e <main>:
   e:   55                      push   %rbp
   f:   48 89 e5                mov    %rsp,%rbp
  12:   48 83 ec 10             sub    $0x10,%rsp
  16:   e8 00 00 00 00          call   1b <main+0xd>
  1b:   66 0f 7e c0             movd   %xmm0,%eax
  1f:   89 45 fc                mov    %eax,-0x4(%rbp)
  22:   b8 00 00 00 00          mov    $0x0,%eax
  27:   c9                      leave
  28:   c3                      ret
            \end{terminalwindow}

\end{frame}

% ------------------------------------------------------------------
% Слайд 12. Линковка: схема
% ------------------------------------------------------------------

\begin{frame}[fragile]

    \frametitle{Линковка}

    \centering

    \begin{tikzpicture}[node distance=5cm]

        \matrix (objs) [row sep=0.5cm, column sep=0.1cm] {
            \node (bin1) {
                \includegraphics[width=1.5cm]
                    {presentations/images/binary-file.png}
            };
 &
            \node (bin2) {
                \includegraphics[width=1.5cm]
                    {presentations/images/binary-file.png}
            };
 \\
            \node (bin3) {
                \includegraphics[width=1.5cm]
                    {presentations/images/binary-file.png}
            };
 &
            \node (bin4) {
                \includegraphics[width=1.5cm]
                    {presentations/images/binary-file.png}
            };
 \\
        };
        \node [below=0mm of bin1] { \texttt{*.o} };
        \node [below=0mm of bin2] { \texttt{*.o} };
        \node [below=0mm of bin3] { \texttt{*.o} };
        \node [below=0mm of bin4] { \texttt{*.o} };

        \node (exe) [right=of objs] {
            \includegraphics[width=1.5cm]
                {presentations/images/binary-file.png}
        };
        \node [below=0mm of exe] { \texttt{ELF} };


        \draw[->, ultra thick] (objs.east) -- (exe.west)
            node[midway, above] {Линковка};

    \end{tikzpicture}

\end{frame}

% ------------------------------------------------------------------
% Слайд 13. Линковка: пример
% ------------------------------------------------------------------

\begin{frame}[fragile]

    \frametitle{Линковка}

    \begin{columns}[T]

        \begin{column}{0.4\textwidth}

            \myinputlisting[minted language=c]
                {presentations/05-translation-phases/src/}
                {link-example.h}

            \myinputlisting[minted language=c]
                {presentations/05-translation-phases/src/}
                {link-example.c}

            \myinputlisting[minted language=c]
                {presentations/05-translation-phases/src/}
                {link-example-main.c}

        \end{column}

        \begin{column}{0.6\textwidth}

            \begin{terminalwindow}[||]
|\shellcommand{gcc -c link-example.c}|
|\shellcommand{gcc -c link-example-main.c}|
|\shellcommand{gcc link-example.o link-example-main.o}|
|\shellcommand{./a.out}|
Hello!
|\shellcommand{file a.out}|
a.out: ELF |\myhighlight{green}{64-bit}| LSB pie executable,
|\myhighlight{green}{x86-64}|, version 1 (SYSV), dynamically linked,
interpreter /nix/store/|\dots|/lib/ld-linux-x86-64.so.2,
for GNU/Linux 3.10.0, not stripped
            \end{terminalwindow}

        \end{column}

    \end{columns}

\end{frame}

% ------------------------------------------------------------------
% Слайд 14. Статическая библиотека: создание архива
% ------------------------------------------------------------------

\begin{frame}[fragile]

    \frametitle{Статическая библиотека}

    \centering

    \begin{tikzpicture}[node distance=5cm]

        \matrix (objs) [row sep=0.5cm, column sep=0.1cm] {
            \node (bin1) {
                \includegraphics[width=1.5cm]
                    {presentations/images/binary-file.png}
            };
 &
            \node (bin2) {
                \includegraphics[width=1.5cm]
                    {presentations/images/binary-file.png}
            };
 \\
            \node (bin3) {
                \includegraphics[width=1.5cm]
                    {presentations/images/binary-file.png}
            };
 &
            \node (bin4) {
                \includegraphics[width=1.5cm]
                    {presentations/images/binary-file.png}
            };
 \\
        };
        \node [below=0mm of bin1] { \texttt{*.o} };
        \node [below=0mm of bin2] { \texttt{*.o} };
        \node [below=0mm of bin3] { \texttt{*.o} };
        \node [below=0mm of bin4] { \texttt{*.o} };

        \node (exe) [right=of objs] {
            \includegraphics[width=1.5cm]
                {presentations/images/binary-file.png}
        };
        \node [below=0mm of exe] { \texttt{lib*.a} };


        \draw[->, ultra thick] (objs.east) -- (exe.west)
            node[midway, above] {Создание архива};

    \end{tikzpicture}

\end{frame}

% ------------------------------------------------------------------
% Слайд 15. Статическая библиотека: линковка
% ------------------------------------------------------------------

\begin{frame}[fragile]

    \frametitle{Статическая библиотека}

    \centering

    \begin{tikzpicture}[node distance=5cm]

        \matrix (objs) [row sep=0.5cm, column sep=0.1cm] {
            \node (bin1) {
                \includegraphics[width=1.5cm]
                    {presentations/images/binary-file.png}
            }; &
            \node (bin2) {
                \includegraphics[width=1.5cm]
                    {presentations/images/binary-file.png}
            }; \\
            \node (bin3) {
                \includegraphics[width=1.5cm]
                    {presentations/images/binary-file.png}
            }; &
            \node (lib) {
                \includegraphics[width=1.5cm]
                    {presentations/images/binary-file.png}
            }; \\
        };
        \node [below=0mm of bin1] { \texttt{*.o} };
        \node [below=0mm of bin2] { \texttt{*.o} };
        \node [below=0mm of bin3] { \texttt{*.o} };
        \node [below=0mm of lib]  { \texttt{lib*.a} };

        \node (exe) [right=of objs] {
            \includegraphics[width=1.5cm]
                {presentations/images/binary-file.png}
        };
        \node [below=0mm of exe] { \texttt{ELF} };

        \draw[->, ultra thick] (objs.east) -- (exe.west)
            node[midway, above] {Линковка};

    \end{tikzpicture}

\end{frame}

% ------------------------------------------------------------------
% Слайд 16. Статическая библиотека: пример
% ------------------------------------------------------------------

\begin{frame}[fragile]

    \frametitle{Статическая библиотека}

    \begin{columns}[T]

        \begin{column}{0.5\textwidth}

            \myinputlisting[minted language=c]
                {presentations/05-translation-phases/src/library/}
                {main.c}

            \myinputlisting[minted language=c]
                {presentations/05-translation-phases/src/library/}
                {foo.c}

            \myinputlisting[minted language=c]
                {presentations/05-translation-phases/src/library/}
                {bar.c}

        \end{column}

        \begin{column}{0.5\textwidth}

            \begin{terminalwindow}[||]
|\shellcommand{gcc -c bar.c -o bar.o}|
|\shellcommand{gcc -c foo.c -o foo.o}|
|\shellcommand{\myhighlight{green}{ar -r libfoo.a} foo.o bar.o}|
ar: creating libfoo.a
|\shellcommand{gcc -c main.c -o main.o}|
|\shellcommand{gcc main.o \myhighlight{green}{-L\$(pwd) -lfoo} -o main}|
|\shellcommand{./main}|
Hello, world!
            \end{terminalwindow}

        \end{column}

    \end{columns}

\end{frame}

% ------------------------------------------------------------------
% Слайд 17. Динамическая библиотека: создание
% ------------------------------------------------------------------

\begin{frame}[fragile]

    \frametitle{Динамическая библиотека}

    \centering

    \begin{tikzpicture}[node distance=5cm]

        \matrix (objs) [row sep=0.5cm, column sep=0.1cm] {
            \node (bin1) {
                \includegraphics[width=1.5cm]
                    {presentations/images/binary-file.png}
            };
 &
            \node (bin2) {
                \includegraphics[width=1.5cm]
                    {presentations/images/binary-file.png}
            };
 \\
            \node (bin3) {
                \includegraphics[width=1.5cm]
                    {presentations/images/binary-file.png}
            };
 &
            \node (bin4) {
                \includegraphics[width=1.5cm]
                    {presentations/images/binary-file.png}
            };
 \\
        };
        \node [below=0mm of bin1] { \texttt{*.o} };
        \node [below=0mm of bin2] { \texttt{*.o} };
        \node [below=0mm of bin3] { \texttt{*.o} };
        \node [below=0mm of bin4] { \texttt{*.o} };

        \node (exe) [right=of objs] {
            \includegraphics[width=1.5cm]
                {presentations/images/binary-file.png}
        };
        \node [below=0mm of exe] { \texttt{lib*.so} };


        \draw[->, ultra thick] (objs.east) -- (exe.west)
            node[midway, above, align=center]
                {Создание динамической\\библиотеки};

    \end{tikzpicture}

\end{frame}

% ------------------------------------------------------------------
% Слайд 18. Динамическая библиотека: линковка
% ------------------------------------------------------------------

\begin{frame}[fragile]

    \frametitle{Динамическая библиотека}

    \centering

    \begin{tikzpicture}[node distance=5cm]

        \matrix (src) [row sep=0.5cm, column sep=0.1cm] {
            \node (bin1) {
                \includegraphics[width=1.5cm]
                    {presentations/images/binary-file.png}
            }; &
            \node (bin2) {
                \includegraphics[width=1.5cm]
                    {presentations/images/binary-file.png}
            }; \\
            \node (bin3) {
                \includegraphics[width=1.5cm]
                    {presentations/images/binary-file.png}
            }; &
            \node (lib0) {
                \includegraphics[width=1.5cm]
                    {presentations/images/binary-file.png}
            }; \\
        };
        \node [below=0mm of bin1] { \texttt{*.o} };
        \node [below=0mm of bin2] { \texttt{*.o} };
        \node [below=0mm of bin3] { \texttt{*.o} };
        \node [below=0mm of lib0] { \texttt{lib*.so} };

        \matrix (dst) [row sep=0.5cm, column sep=0.1cm, right=of src] {
            \node (lib1) {
                \includegraphics[width=1.5cm]
                    {presentations/images/binary-file.png}
            }; \\
            \node (exe) {
                \includegraphics[width=1.5cm]
                    {presentations/images/binary-file.png}
            }; \\
        };
        \node [below=0mm of exe] { \texttt{executable} };
        \node [below=0mm of lib1] { \texttt{lib*.so} };

        \draw[->, ultra thick] (src.east) -- (dst.west)
            node[midway, above, align=center] {Линковка};

    \end{tikzpicture}

\end{frame}

% ------------------------------------------------------------------
% Слайд 19. Динамическая библиотека: пример
% ------------------------------------------------------------------

\begin{frame}[fragile]

    \frametitle{Динамическая библиотека}

    \begin{columns}[T]

        \begin{column}{0.4\textwidth}

            \myinputlisting[minted language=c]
                {presentations/05-translation-phases/src/library/}
                {main.c}

            \myinputlisting[minted language=c]
                {presentations/05-translation-phases/src/library/}
                {foo.c}

            \myinputlisting[minted language=c]
                {presentations/05-translation-phases/src/library/}
                {bar.c}

        \end{column}

        \begin{column}{0.6\textwidth}

            \begin{terminalwindow}[||]
|\shellcommand{gcc -c -fPIC bar.c -o bar.o}|
|\shellcommand{gcc -c -fPIC foo.c -o foo.o}|
|\shellcommand{gcc \myhighlight{green}{-shared} bar.o foo.o -o libfoo.so}|
|\shellcommand{gcc -c main.c -o main.o}|
|\shellcommand{gcc main.o \myhighlight{green}{-L\$(pwd) -lfoo -Wl,-rpath=\$(pwd)} -o main}|
|\shellcommand{./main}|
Hello, world!
            \end{terminalwindow}

        \end{column}

    \end{columns}

\end{frame}

% ------------------------------------------------------------------
% Слайд 20. Релокации: постановка проблемы
% ------------------------------------------------------------------

\begin{frame}[fragile]

    \frametitle{Релокации}

    Во время компиляции одной единицы трансляции компилятор
    не может знать адреса функций, которые не были определены
    в данной единице трансляции.

    \bigbreak

    % In-place listing: an educational fragment without a file
    % name, with an explanatory comment; the same code without
    % the comment is shown on slide 23 (relocations/bar.c).
    \begin{myinplacelisting}[minted language=c]
void foo(void);

void bar(void) {
    // Which address to put in assembly for `foo`?
    foo();
}
    \end{myinplacelisting}

\end{frame}

% ------------------------------------------------------------------
% Слайд 21. Релокации: что делает компилятор
% ------------------------------------------------------------------

\begin{frame}

    \frametitle{Релокации}

    На этапе компиляции единицы трансляции компилятор сделает
    две вещи:

    \begin{enumerate}

        \item В инструкцию \texttt{call} компилятор вместо адреса
            \texttt{foo} положит нули.

        \item Компилятор создаст \textit{таблицу релокации}
            в объектном файле, в которой поставит в соответствие
            адресу первого нуля имя нужной функции.

    \end{enumerate}

\end{frame}

% ------------------------------------------------------------------
% Слайд 22. Релокации: objdump -dC и -rC
% ------------------------------------------------------------------

\begin{frame}[fragile]

    \frametitle{Релокации}

            \begin{terminalwindow}
!\shellcommand{gcc -c bar.c}!

!\shellcommand{objdump -dC bar.o}!

0000000000000000 <bar>:
   0:   55                      push   %rbp
   1:   48 89 e5                mov    %rsp,%rbp
   !\myhighlight{green}{4}!:   e8 !\myhighlight{yellow}{00 00 00 00}!          call   9 <bar+0x9>
   9:   90                      nop
   a:   5d                      pop    %rbp
   b:   c3                      ret

!\shellcommand{objdump -rC bar.o}!

RELOCATION RECORDS FOR [.text]:
OFFSET           TYPE              VALUE
00000000000000!\myhighlight{green}{05}! R_X86_64_PLT32    foo-0x0000000000000004
            \end{terminalwindow}

    \bigbreak

    Первый \texttt{00} байт в инструкции \texttt{call} имеет адрес
    \texttt{5} (\texttt{4}-й адрес занимает opcode инструкции).

\end{frame}

% ------------------------------------------------------------------
% Слайд 23. Релокации: работа линковщика
% ------------------------------------------------------------------

\begin{frame}[fragile]

    \frametitle{Релокации}

    На этапе линковки линковщик, используя таблицы релокаций,
    заполнит все нулевые адреса.

    \begin{columns}[T]

        \begin{column}{0.4\textwidth}

            \myinputlisting[minted language=c]
                {presentations/05-translation-phases/src/relocations/}
                {bar.c}

            \myinputlisting[minted language=c]
                {presentations/05-translation-phases/src/relocations/}
                {foo.c}

        \end{column}

        \begin{column}{0.6\textwidth}

            \begin{terminalwindow}
!\shellcommand{gcc bar.c foo.c -o main}!

!\shellcommand{objdump -dC main}!

0000000000001129 <bar>:
    1129:       55
    112a:       48 89 e5
    112d:       e8 !\myhighlight{yellow}{03 00 00 00}!
    !\dots!

!\dots!

0000000000001143 <main>:
    1143:       55
    1144:       48 89 e5
    1147:       e8 !\myhighlight{yellow}{dd ff ff ff}!
    !\dots!

!\dots!
            \end{terminalwindow}

        \end{column}

    \end{columns}

\end{frame}

% ------------------------------------------------------------------
% Слайд 24. Релокации: формула адреса
% ------------------------------------------------------------------

\begin{frame}[fragile]

    \frametitle{Релокации}

    Адрес функции из операнда \texttt{call} считается по формуле:

    \begin{center}
    \textit{address} = \textit{next} + \textit{operand}
    \end{center}

    где \textit{next} --- адрес инструкции, следующей за
    \texttt{call} (\texttt{0x1132}); \textit{operand} --- смещение,
    записанное в операнде \texttt{call} (\texttt{0x03}).

    \begin{center}
    \texttt{0x1132} + \texttt{0x03} = \texttt{0x1135}
    \end{center}

    \bigbreak

            \begin{terminalwindow}
0000000000001129 <bar>:
    112d:       e8 !\myhighlight{yellow}{03 00 00 00}!          call   1135 <foo>
    !\myhighlight{yellow}{1132}!:       90                      nop

000000000000!\myhighlight{green}{1135}! <foo>:
    !\myhighlight{green}{1135}!:       55                      push   %rbp
            \end{terminalwindow}

\end{frame}

% ------------------------------------------------------------------
% Слайд 25. Источники, документация и литература
% ------------------------------------------------------------------

\begin{frame}

    \frametitle{Источники, документация \\ и литература}

    \begin{itemize}

        \item Статья про фазы трансляции C на cppreference:
            \url{https://en.cppreference.com/w/c/language/translation_phases}

        \item Перевод статьи Beginner's guide to linkers на Хабре:
            \url{https://habr.com/ru/articles/150327/}

        \item Цикл статей про архитектуру компьютера с фокусом
            на выполнение процессов: \url{https://cpu.land/}

        \item Создание статических и динамических библиотек
            при помощи \texttt{gcc}:
    \url{https://renenyffenegger.ch/notes/development/languages/C-C-plus-plus/GCC/create-libraries}

    \end{itemize}

\end{frame}

\end{document}
